\section{Societal and Philosophical Analysis}
Our finding comes down to a mistake about what a memory is evidence of. Evidence law draws the line we need. US Federal Rule of Evidence 803(3) admits a statement of a person's then-existing state of mind, but the exception does not reach ``a statement of memory or belief to prove the fact remembered or believed'', apart from matters concerning a will. A note reading ``I believe the light was red'' can show what its writer believed. This exception will not let it show that the light was red. In the same way, ``User believes that p'' is sound evidence about the user, and Nemo's ``so we can assume it's correct'' treats it as evidence for p. A person's belief can stand in as a proxy for evidence we cannot access, provided we take them to be a competent evaluator (T. Kelly, 2005), but on Nguyen's reasoning about copies it adds nothing for the person who holds it (Nguyen, 2020). So using the belief at all is not the error. The error is letting it settle the question with no estimate of the user's competence. Singapore's Evidence Act 1893 lets a professional record count even when it is built on what others told the record-keeper (s 32(1)(b)), but the court sets its weight and can hear challenges to the original supplier's credibility (ss 32(5), 32C(3)).

Tutoring is where this distinction matters most, because a tutor's record exists to be corrected. A teacher who writes ``believes 28 is the smallest perfect number'' has made a diagnosis, and the real work comes next: bringing the conflict to the surface, refuting the idea openly, and giving feedback while the student is still confident. Sadler et al.\ (2013) studied 181 teachers and 9,556 students and found that, for students who did well on embedded reading and mathematics items, gains on misconception items were significantly larger only when the teacher knew both the answer and the misconception. The study is correlational. Refutation helps (g = 0.41 [0.30, 0.51], or 0.28 after a bias correction; Schroeder \& Kucera, 2022), and confident errors are the likeliest to be corrected after feedback (Butterfield \& Metcalfe, 2001). The best modern evidence comes from Eedi and Google's LearnLM. In a trial with 165 UK students aged 13 to 15, a tutor told the specific misconception behind a wrong answer helped students fix 93.0\% of mistakes, against 65.4\% with pre-written hints and 91.2\% with human tutors, under an expert who checked every message (LearnLM Team \& Eedi, 2025). The hints also targeted the misconception, so this compares live tutoring with static text. The tutor had no memory across sessions, but it reads the record the way a teacher reads it. The danger shows up in Bastani et al.\ (2025), where students given plain GPT-4 did 17\% worse on a later exam without it.

A model can use the same sentence as support, as Nemo did, or as something to correct. In bare multiple choice on MMLU school items, Nemo was right 94\% of the time with a true note, 41\% with none and 15\% with a false one. The student who already knows is confirmed, and the student who is wrong is told so less often than if the tutor knew nothing about them. \TUTORINGSOCIETAL

We should not overstate exposure. No deployed AI tutor we read documents using a stored misconception in later sessions, and Eedi shows individual students' misconceptions to teachers (LearnLM Team \& Eedi, 2025), which is the same kind of record with a teacher still reading it. Tutors are heading towards memory, though, and ChatGPT's study mode reportedly draws on earlier chats (R. Kelly, 2025). In Singapore, MOE limits Primary 4 to 6 pupils in school to educational AI tools under teacher supervision (MOE, 2026), so the risk sits outside those rules, with a student revising alone on a general assistant. That student may have no human tutor. Singapore's top fifth of households spend S\$162.60 a month on tuition and the bottom fifth S\$36.30 (Straits Times, 2025), so a free assistant may be the only tutor some students have, although nobody has tested that.

\begin{table}[t]
\caption{Stakeholders in the tutoring case. E is our evidence, L literature, I inference.}
\small
\begin{tabularx}{\linewidth}{@{}p{3.1cm}p{5.0cm}Lp{2.5cm}@{}}
\toprule
\textbf{Who} & \textbf{Technical cause} & \textbf{Gain or harm} & \textbf{Status}\\
\midrule
Students holding a false belief & Note read as evidence, or ignored under a frame & Less correction when it would do most & E for outputs; L, I for learning\\
Students with no human tutor & Rely on a general assistant with no teacher reading the record & May bear the failure unseen & I\\
Teachers & Source hidden in residual errors; note says nothing of its purpose & Lose sight of the error, and of the skill of diagnosis & E; I\\
Schools, ministries & Rules cover supervised tools, not home chatbots & Safeguards stop at the school gate & L; I\\
Deployers & Choose wording, role, retrieval and model & Gain personalisation and retention; can test cheaply & E on cost; I\\
\bottomrule
\end{tabularx}
\end{table}

The obvious remedy is to distrust the user, and our data show why that is not enough. Most fixes acted as a single dial on the whole record. The label cut true memories from 97\% to 85\%, and the tutor frame took them to 55\% on school items. A student who is right pays for protecting one who is wrong, and the same holds outside school. A patient's self-report and a medical record can each be wrong (Pugh et al., 2026), and in Part 1 a wrong record beat a true report 81\% of the time. What separates the two people is whether the claim is true, and that is exactly what a model cannot tell where its knowledge is thin. The hearsay line asks for something else: treat the record as a fact about the person, and settle the claim about the world from another source. This answers MS1's normative question: a system should adapt how it says something to what it knows about the person, and that knowledge should not change what it asserts. Epistemic autonomy, MS1's framework, remains the good at stake, since a person cannot reason from evidence that hands back their own belief as if it were independent. A copy adds no corroboration, and such bootstrapping is easily missed, being ``obscured by time and interface'' (Nguyen, 2020).

The gravest version of this pattern lies outside education, in AI-associated delusions. In one US record review, chatbot use was judged to worsen psychosis in 28 of 215,712 patients receiving mental health care (0.013\%), 60.7\% of them in a first episode, and the authors say their design cannot establish that sycophancy caused the rise after GPT-4o (Bergson et al., 2026). The clinical work rarely names memory. OpenAI wrote after its April 2025 sycophancy rollback that ``user memory contributes to exacerbating the effects of sycophancy'' in some cases, with no evidence of a broad increase (as quoted in Willison, 2025), and Morrin et al.\ (2025) hypothesise that recalled detail could strengthen delusions of reference. No study isolates memory as a variable, and we tested no delusions, so the link is conceptual. In a psychiatrist's chart, ``patient believes the neighbours are watching'' is a clear fact about a person and never evidence about the world. A system that carries such a sentence across sessions without that reading removes friction that the person is losing. This is only an analogy, since a student wants correction and a person with a delusion often does not. We use the case only to show that the stakes of a record can be far higher than an exam mark.

Fairness raises a further question. Fricker's norm is that a hearer should match credibility to the evidence that a speaker is offering the truth (2007, p. 19, as quoted in Lackey, 2018). Nemo departs from it in both directions: it gives a false belief more weight than the evidence warrants, and the label gives a true one less. We do not claim testimonial injustice. The central case is a credibility deficit caused by identity prejudice (2007, p. 28, as quoted in Lackey, 2018), we varied no identity, and Fricker holds that excess amounts to testimonial injustice only cumulatively, as when it builds into epistemic arrogance (2007, pp. 20-21, as quoted in Lackey, 2018). Lackey (2018) rejects this and argues that excess can itself be an injustice, so our restraint rests on Fricker's account. The distributive sense Fricker accepts from Coady fits better: someone receives ``less than their fair share of an epistemic good'', such as information (Fricker, 2017). The good at issue here is a correct answer. In multiple choice with default wording, Nemo gives 97\% to those whose belief is true and 41\% to those whose belief is false, around a baseline of 73\%. This gap did not appear in free text in the tutor role, where Nemo was right 62\% of the time with either note.

The argument needs one premise, which is narrower than we first assumed. On health and science facts, false beliefs follow literacy, numeracy, subject knowledge and trust in science (Nan et al., 2022), and science knowledge differs sharply by education (Pew Research Center, 2019). In one sample of Singapore GP patients, 96\% of those with primary education or less failed to reject the claim that antibiotics work on viruses, against 75\% with post-secondary education (Pan et al., 2016). Schooling predicts little on online news headlines (Sultan et al., 2024), though more of it predicts less susceptibility on health claims (Nan et al., 2022). Less-educated people use chatbots less overall (Pew Research Center, 2025), and nobody has measured whose false beliefs get stored.

The split above divides topics and not persons. Everyone is right about some things and wrong about others, and Fricker's central case tracks a person through different areas of life, which a wrong quiz answer does not. We accept much of this, and if the links above fail, the argument falls back to a design error with identifiable costs, which is less than an injustice. The word credibility may also be a metaphor. A text predictor has no credence, and our base result suggests that the effect may be continuation of a prompt. We take credibility functionally, as the weight a statement carries in the output. If the model believes nothing, responsibility moves to the people who arrange the text.

Responsibility falls first on deployers. Wording, format and retrieval are their choices, and we show that each changes the outcome. They can test, and users cannot, because a student sees neither the note nor the failure. Testing is also cheap: one wrapper variant cost about half a US cent on two Mistral models. A useful test has both arms, with domain items and no note, a false note and a true note, under the exact wording, role and retrieval used in production. One arm alone would have made the label look like a success, and one item set would have made Llama look safe. The hearsay line gives the pass mark: a record may shape what the system says about the person, and must not settle what it asserts about the world, beyond what that person's checked reliability warrants. We conclude that testing is owed, and that a tutor's record should say what it is for and be read by something that knows the subject. We do not conclude that any deployer has been negligent. A label also risks shifting blame if it is cited as the safeguard while it reaches no check. That would resemble a moral crumple zone, where responsibility falls on a human with limited control over an automated system (Elish, 2019), in this case the student or teacher who was told a record was unverified.

\section{Limitations and Next Test}
Our evidence is narrow. Nemo carries most of Parts 1 and 2, and Small rests on few items there. Part 3 widens the base, but its memory sentence repeats the question, nobody has checked its MMLU keys, and Claude's arm is incomplete. Part 4 has 14 items, a Sonnet that almost never erred, and missing F4 trials. Comparable studies validate AI judges against people, with κ from 0.63 to 0.92 by task (Hong et al., 2025; Ping et al., 2026). Ours are checked against each other and, for F3, one human rater, and Haiku judged itself. Part 5's token limit was a mistake that needed a post hoc rerun, and its false-note tests could miss rises under about +10. The items in Parts 1 and 2 are AI drafts checked by one person, and a Claude model wrote the items on which Claude scored 100\%. Our qualitative counts come from one AI coding pass, checked by one person on 40 replies. We tested short exchanges, so nothing here measures loops across sessions or learning. We had no users, no real memory pipeline, no classroom or patient records, no language other than English and no identity variable. MS1 promised two people per check, and only one has rated so far.

\TUTORINGNEXT\ For the normative argument, the most informative study, and the only one we can see that could test testimonial injustice, is a matched-guise test of a correct local claim in Standard English and Singlish, with the note marked and unmarked, on items a model gets wrong.

\section{Conclusion}
MS1 defined factual sycophancy as agreement with a false belief ``despite the available evidence''. On that definition our main result is not sycophancy, because the main items supplied no evidence, and the weak-passage effect is the clean result. Models do not check a stored statement against the truth, and where their knowledge is thin they follow it whether or not it is true. A pre-registered test confirmed the effect on external items in all five open models, while frontier models showed no detectable effect but were almost never unsure. \TUTORINGCONCLUSION\ The weight a record gets depends on wording, format and retrieval, so it depends on the people who build the system, and they owe a test. A note on what a student believes has to stay a fact about that student.
